%% toc: P15 164 164.10
\section{Proofs for uniqueness, Cox's postulates and non-additive belief}
\label{sec:annex-p15-w02b-p15-032}

Throughout, \(S\) is finite and nonempty, \(\Delta(S)\) is the set of
probability vectors on \(S\), and \(E_P(f)=\sum_{s\in S}P(s)f(s)\).

\begin{annexitem}{annex:p15-w02b-p15-032-unique}{thm:p15-w02b-p15-032-unique}{Fair prices, symmetry and the vacuous schedule}
\begin{proof}
For part~1, a two-sided schedule is exactly a quote schedule of
\cref{def:p17-finite-price-coherence}: a negative coefficient on
\(b_A=\mathbf 1_A-q(A)\mathbf 1\) is a sale at the quoted price. By
\cref{thm:p17-finite-price-coherence}, avoiding a strict sure loss makes
\(q\) a probability measure, given by a unique vector \(\pi\). A compatible
\(P\) must satisfy \(P(A)\geq q(A)\), because the agent buys at \(q(A)\),
and \(P(A)\leq q(A)\), because the agent sells at \(q(A)\). Hence
\(P(A)=q(A)\) for every \(A\). In particular \(P(\{s\})=\pi_s\) for each
\(s\), so \(P=\pi\). Conversely \(\pi\) satisfies every constraint with
equality. The compatible set is therefore \(\{\pi\}\), which is \(\{q\}\).

For part~2, let every singleton price equal \(c\). By additivity,
\(1=q(S)=\sum_{s\in S}q(\{s\})=|S|c\), so \(c=1/|S|\), and \(q\) is uniform.

For part~3, a one-sided portfolio buys tickets only, so it is
\(\sum_ja_j(\mathbf 1_{A_j}-q(A_j)\mathbf 1)\) with every \(a_j\geq0\).
A term with \(A_j\neq S\) has price zero and nonnegative payoff in every
state. A term with \(A_j=S\) pays \(1-1=0\) in every state. Every such
portfolio is therefore nonnegative in every state, and no strict sure loss
exists. The compatible set consists of the \(P\) with \(P(S)\geq1\) and
\(P(A)\geq0\) for every \(A\), which is all of \(\Delta(S)\). Finally, if
\(A\neq S\), choose \(s\notin A\); the point mass at \(s\) gives \(A\)
probability zero, so \(\inf_{P\in\Delta(S)}P(A)=0=q(A)\). Also
\(\inf_PP(S)=1=q(S)\). Every price equals the lower envelope of the
compatible set.
\end{proof}\end{annexitem}

\begin{annexitem}{annex:p15-w02b-p15-032-cox}{thm:p15-w02b-p15-032-cox}{Product rule, forced endpoints and the ignorance exclusion}
\begin{proof}
For part~1, suppose \(P(C\cap B)>0\). Then \(P(B)\geq P(C\cap B)>0\), so
both conditionals on the right are defined, and
\[
 \frac{P(A\cap C\cap B)}{P(B)}
 =\frac{P(A\cap C\cap B)}{P(C\cap B)}\cdot\frac{P(C\cap B)}{P(B)}.
\]
The left side is \(g(A\cap C\mid B)\), and the right side is
\(F(g(A\mid C\cap B),g(C\mid B))\) with \(F(x,y)=xy\). If \(P(B)>0\), then
\(P(A^c\cap B)=P(B)-P(A\cap B)\), and dividing by \(P(B)\) gives
\(g(A^c\mid B)=1-g(A\mid B)\). That is postulate~(ii) with \(n(x)=1-x\).

For part~2, fix \(B\) and apply postulate~(ii) with \(A=S\). Since
\(S^c=\varnothing\), it gives \(0=g(\varnothing\mid B)=n(g(S\mid B))=n(1)\).
Applying it with \(A=\varnothing\) gives \(1=g(S\mid B)=n(g(\varnothing\mid B))=n(0)\).
The same function \(n\) serves for every event. So if
\(g(A\mid B)=0\), then \(g(A^c\mid B)=n(0)=1\). Hence \(g(A\mid B)\) and
\(g(A^c\mid B)\) cannot both be zero.

For part~3, suppose some \(n\) satisfied \(N(A^c)=n(N(A))\) for all \(A\).
Since \(N(\varnothing)=0\) and \(N(S)=1\), part~2 applied with \(B=S\) gives
\(n(0)=1\). The given event has \(N(A)=0\), so \(N(A^c)=1\), contradicting
\(N(A^c)=0\). For the example, \(\Pi(A)=\max_{s\in A}\pi(s)\) and
\(N(A)=1-\Pi(A^c)\). Thus \(N(\{1\})=1-\Pi(\{2,3\})=0\),
\(N(\{2,3\})=1-\Pi(\{1\})=0\), \(N(\{3\})=1-\Pi(\{1,2\})=0\) and
\(N(\{1,2\})=1-\Pi(\{3\})=1\). Here \(\{1\}\) and \(\{3\}\) have equal
necessity while their complements do not, so no function \(n\) can exist.
\end{proof}\end{annexitem}

\begin{annexitem}{annex:p15-w02b-p15-032-disjoint}{thm:p15-w02b-p15-032-disjoint}{Infimum of a sum and the envelope bounds}
\begin{proof}
Every \(P\in\mathcal P\) has \(P(\varnothing)=0\) and \(P(S)=1\), so the
infima are \(0\) and \(1\). If \(A\subseteq B\), then \(P(A)\leq P(B)\) for
each \(P\), and taking infima gives \(P_*(A)\leq P_*(B)\). For disjoint
\(A,B\) and each \(P\in\mathcal P\),
\[
 P(A\cup B)=P(A)+P(B)\geq P_*(A)+P_*(B).
\]
The right side is a lower bound for every \(P(A\cup B)\), so it is at most
their infimum \(P_*(A\cup B)\). Finally \(P_*(A)\leq P(A)\) by definition of
the infimum. Applied to \(A^c\), the same fact gives \(P_*(A^c)\leq P(A^c)\),
hence \(P(A)=1-P(A^c)\leq1-P_*(A^c)=P^*(A)\) for every \(P\in\mathcal P\).
\end{proof}\end{annexitem}

\begin{annexitem}{annex:p15-w02b-p15-032-ellsberg}{thm:p15-w02b-p15-032-ellsberg}{Superadditivity of lower expectations and the urn values}
\begin{proof}
All gambles are bounded on finite \(S\), so every \(LP(f)\) is a finite real
number. For any gambles \(u,v\) and any \(P\in\mathcal P\),
\(E_P(u+v)=E_P(u)+E_P(v)\geq LP(u)+LP(v)\), and taking the infimum gives
\(LP(u+v)\geq LP(u)+LP(v)\).

For part~1, apply this with \(u=f-g\) and \(v=g\):
\(LP(f)\geq LP(f-g)+LP(g)\). If \(LP(f-g)\geq0\), then \(LP(f)\geq LP(g)\).
If \(LP(f-g)>0\), then \(LP(f)>LP(g)\). Reflexivity holds because
\(LP(0)=0\). For transitivity, if \(LP(f-g)\geq0\) and \(LP(g-h)\geq0\), then
\(LP(f-h)\geq LP(f-g)+LP(g-h)\geq0\).

For part~2, an infimum of real numbers is nonnegative exactly when every
one of the numbers is nonnegative. By linearity \(E_P(f-g)=E_P(f)-E_P(g)\),
so \(LP(f-g)\geq0\) holds exactly when \(E_P(f)\geq E_P(g)\) for every
\(P\in\mathcal P\). No attainment of the infimum was used.

For part~3, at \(s\in E\) the difference \(f_Eh-g_Eh\) is \(f(s)-g(s)\), and
at \(s\notin E\) it is \(h(s)-h(s)=0\). Thus \(f_Eh-g_Eh=(f-g)\mathbf 1_E\),
which does not depend on \(h\). Both \(\succeq_D\) and \(\succ_D^{+}\) are
defined from \(LP\) of this difference alone. Each comparison therefore
has the same truth value for \(h\) and for \(h'\).

For part~4, in the order \((R,B,Y)\) the acts are \(F_1=(1,0,0)\),
\(F_2=(0,1,0)\), \(F_3=(1,0,1)\) and \(F_4=(0,1,1)\). For
\(P=(1/3,b,2/3-b)\), the expectations are \(1/3\), \(b\), \(1-b\) and
\(2/3\). Over \(0\leq b\leq2/3\) their minima are \(1/3\), \(0\) (at \(b=0\)),
\(1/3\) (at \(b=2/3\)) and \(2/3\). Hence \(F_1\succ_MF_2\) and
\(F_4\succ_MF_3\). Take \(E=\{R,B\}\), let \(h\) be zero and let \(h'\) be one.
Then \(F_1=(F_1)_Eh\), \(F_2=(F_2)_Eh\), \(F_3=(F_1)_Eh'\) and
\(F_4=(F_2)_Eh'\). The pasting condition would require
\(F_1\succeq_MF_2\) exactly when \(F_3\succeq_MF_4\). The first holds and
the second fails, since \(1/3<2/3\). For difference comparison,
\(F_1-F_2=(1,-1,0)=F_3-F_4\), with expectation \(1/3-b\), ranging over
\([-1/3,1/3]\). So \(LP(F_1-F_2)=-1/3<0\), and likewise
\(LP(F_2-F_1)=\min_b(b-1/3)=-1/3<0\). Neither of \(F_1,F_2\) is weakly
preferred to the other under \(\succeq_D\). The pair \(F_3,F_4\) has the
same differences and hence the same verdict.

For part~5, the belief function is
\(\operatorname{Bel}(A)=\sum_{\varnothing\neq E\subseteq A}m(E)\). It
gives \(1/3\) to every event containing \(R\) but not both \(B\) and \(Y\),
\(2/3\) to \(\{B,Y\}\), \(1\) to \(S\), and \(0\) to \(\{B\}\), \(\{Y\}\) and
\(\varnothing\). A probability \(P\) dominates it exactly when
\(P(\{R\})\geq1/3\) and \(P(\{B,Y\})\geq2/3\), since every other constraint
follows from these by monotonicity. As the two values sum to one, both hold
with equality. These are exactly the members of \(\mathcal P_E\). The
pignistic formula of \cref{thm:p15-r47-pignistic} gives
\(\operatorname{BetP}(R)=1/3\) and
\(\operatorname{BetP}(B)=\operatorname{BetP}(Y)=(2/3)/2=1/3\). Under this
uniform probability the expected utilities of \(F_1,F_2,F_3,F_4\) are
\(1/3,1/3,2/3,2/3\), so \(F_1\sim F_2\) and \(F_3\sim F_4\). For \(m(S)=1\),
the formula gives \(\operatorname{BetP}(s)=1/|S|\) for each \(s\).
\end{proof}\end{annexitem}

\begin{annexitem}{annex:p15-w02b-p15-032-mel}{thm:p15-w02b-p15-032-mel}{Containment semantics for K, D and ignorance}
\begin{proof}
For part~1, \(E\subseteq[p\land q]=[p]\cap[q]\) holds exactly when
\(E\subseteq[p]\) and \(E\subseteq[q]\). This is the K-instance, and in set
terms it is \(N_E(A\cap B)=\min\{N_E(A),N_E(B)\}\). For the schema, suppose
\(E\subseteq[p\Rightarrow q]\) and \(E\subseteq[p]\). Every \(s\in E\) then
satisfies \(p\) and \(p\Rightarrow q\), hence \(q\), so \(E\subseteq[q]\).
Nothing here uses \(E\neq\varnothing\).

For part~2, let \(E\neq\varnothing\). If \(E\subseteq[p]\) and
\(E\subseteq[\neg p]\), then \(E\subseteq[p]\cap[\neg p]=\varnothing\),
which is impossible. So \(\Box p\Rightarrow\neg\Box\neg p\) holds. If instead
\(E=\varnothing\), both \(\Box p\) and \(\Box\neg p\) hold for any \(p\), and
the schema fails. In set terms, \(\Pi_E(A)=1-N_E(A^c)\) equals one exactly
when \(E\not\subseteq A^c\), that is, when \(E\cap A\neq\varnothing\). If
\(E\neq\varnothing\) and \(E\subseteq A\), then \(E\cap A=E\neq\varnothing\).
If \(E=\varnothing\), then \(N_E(\varnothing)=1\) and \(\Pi_E(\varnothing)=0\).
Also \(N_E(\varnothing)=0\) holds exactly when \(E\not\subseteq\varnothing\).

For part~3, \(\neg\Box p\) means \(E\not\subseteq[p]\), that is, \(E\) meets
\([\neg p]\). Likewise \(\neg\Box\neg p\) means that \(E\) meets \([p]\).

For part~4, \cref{thm:p15-r47-threshold} gives
\(\min\{N(A),N(A^c)\}=0\). If \(N(A)>0\), then \(N(A^c)=0\), so
\(\Pi(A)=1-N(A^c)=1\). With \(A=[p]\) and \(N([p])\geq\lambda>0\), this gives
\(\Pi([p])=1\), which is the truth condition of \(\Diamond^1p\).
\end{proof}\end{annexitem}

\begin{annexitem}{annex:p15-w02b-p15-032-chain}{thm:p15-w02b-p15-032-chain}{Order-only representation and contextual acceptance}
\begin{proof}
For part~1, every \(A\subseteq S\) equals
\(\bigcap_{s\notin A}(S\setminus\{s\})\); for \(A=S\) this is the empty
intersection \(S\). Repeated minitivity over this finite intersection gives
\(N(A)=\min_{s\notin A}N(S\setminus\{s\})\), and for \(A=S\) the value is
\(N(S)=1_L\). The argument uses only the order of \(L\). With
\(A=\varnothing\), the minimum over all \(s\) equals \(N(\varnothing)=0_L\).
A minimum over a finite set is attained, so some \(s\) has
\(N(S\setminus\{s\})=0_L\). Now let \(L=\{0,1\}\) and
\(E=\{s:N(S\setminus\{s\})=0\}\), which is nonempty. Then \(N(A)=1\) exactly
when every \(s\notin A\) has \(N(S\setminus\{s\})=1\), that is, when no
\(s\notin A\) lies in \(E\), that is, when \(E\subseteq A\). So \(N=N_E\).
Conversely, for nonempty \(E\), part~1 of \cref{thm:p15-w02b-p15-032-mel}
gives minitivity, \(N_E(S)=1\), and \(N_E(\varnothing)=0\) because
\(E\not\subseteq\varnothing\).

For part~2, \([p\Rightarrow q]^c=[p\land\neg q]\) and
\([p\Rightarrow\neg q]^c=[p\land q]\). With \(N(A)=1-\Pi(A^c)\),
\[
 N([p\Rightarrow q])>N([p\Rightarrow\neg q])
 \iff 1-\Pi([p\land\neg q])>1-\Pi([p\land q])
 \iff \Pi([p\land q])>\Pi([p\land\neg q]).
\]

For part~3, \([p]\) is the disjoint union of \([p\land q]\) and
\([p\land\neg q]\), so \(\Pi([p])\) is the larger of their possibilities.
Thus \(q\) is accepted in context \(p\) exactly when
\(\Pi([p\land q])=\Pi([p])>\Pi([p\land\neg q])\). Suppose \(q_1\) and \(q_2\)
are accepted. Since \([p\land\neg(q_1\land q_2)]\) is the union of
\([p\land\neg q_1]\) and \([p\land\neg q_2]\), its possibility is the larger
of two values, each below \(\Pi([p])\). Since \([p]\) is the union of
\([p\land q_1\land q_2]\) and \([p\land\neg(q_1\land q_2)]\), and the second
part has possibility below \(\Pi([p])\), the first part has possibility
\(\Pi([p])\). So \(q_1\land q_2\) is accepted. Accepting both \(q\) and
\(\neg q\) would require \(\Pi([p\land q])>\Pi([p\land\neg q])\) together
with the reverse strict inequality, which is impossible.

For part~4, the maximum of \(\pi\) is one, so \(\pi\) is normalized. In
context \(b\), the valuations satisfying \(b\land f\) are \(bgf\) and
\(b\bar gf\), with possibility \(\max\{1/4,1\}=1\). Those satisfying
\(b\land\neg f\) are \(bg\bar f\) and \(b\bar g\bar f\), with possibility
\(\max\{1/2,1/4\}=1/2\). As \(1>1/2\), \(f\) is accepted. In context
\(b\land g\), the possibility of \(bgf\) is \(1/4\) and that of \(bg\bar f\)
is \(1/2\). As \(1/2>1/4\), \(\neg f\) is accepted in the stronger context,
and by part~3 \(f\) is not.
\end{proof}\end{annexitem}

\begin{annexitem}{annex:p15-w02b-p15-032-convex}{thm:p15-w02b-p15-032-convex}{Nearest-point separation in the simplex}
\begin{proof}
Identify \(\Delta(S)\) with a compact convex subset of \(\mathbb R^S\), and
write \(\langle f,P\rangle=E_P(f)\). Let \(\bar K\) be the closure of \(K\).
It lies in the closed set \(\Delta(S)\) and is therefore compact. It is
convex, because a limit of convex combinations of points of \(K\) is a
convex combination of their limits. For each \(f\) the map
\(P\mapsto\langle f,P\rangle\) is continuous, so
\(\inf_{\bar K}\langle f,P\rangle=\inf_K\langle f,P\rangle=LP_K(f)\).

First, \(\bar K\subseteq\mathcal M(LP_K)\). If \(P\in\bar K\), then
\(\langle f,P\rangle\geq\inf_{\bar K}\langle f,\cdot\rangle=LP_K(f)\) for
every \(f\).

Second, let \(Q\in\Delta(S)\setminus\bar K\). The continuous function
\(P\mapsto\|P-Q\|^2\) attains its minimum on the compact set \(\bar K\) at
some \(P_0\), and \(P_0\neq Q\). For \(P\in\bar K\) and \(0<t\leq1\), the point
\(P_0+t(P-P_0)\) lies in \(\bar K\). Hence
\[
 0\leq\|P_0+t(P-P_0)-Q\|^2-\|P_0-Q\|^2
   =2t\langle P_0-Q,P-P_0\rangle+t^2\|P-P_0\|^2 .
\]
Dividing by \(t\) and letting \(t\to0\) gives
\(\langle P_0-Q,P-P_0\rangle\geq0\). Put \(f=P_0-Q\). Then for every
\(P\in\bar K\),
\[
 \langle f,P\rangle\geq\langle f,P_0\rangle
 =\langle f,Q\rangle+\|P_0-Q\|^2>\langle f,Q\rangle .
\]
Taking the infimum over \(\bar K\) gives
\(LP_K(f)\geq\langle f,P_0\rangle>E_Q(f)\), so \(Q\notin\mathcal M(LP_K)\).
Therefore \(\mathcal M(LP_K)=\bar K\), and it equals \(K\) exactly when \(K\)
is closed.

In the example, \(\bar K=\{(t,1-t):0\leq t\leq1\}\) contains \((0,1)\), which
is not in \(K\). Directly, for any \(f=(f_1,f_2)\),
\(LP_K(f)=\inf_{0<t\leq1}(tf_1+(1-t)f_2)=\min\{f_1,f_2\}\leq f_2=E_{(0,1)}(f)\).
\end{proof}\end{annexitem}

\begin{annexitem}{annex:p15-w02b-p15-032-big-stepped}{thm:p15-w02b-p15-032-big-stepped}{Dominant states, co-singletons and the exact condition}
\begin{proof}
Fix a context \(C\) with \(T=P(C)>0\). Because
\(P(A\cap C)+P(A^c\cap C)=T\), the condition \(P(A\mid C)>P(A^c\mid C)\)
holds exactly when \(P(A\cap C)>T/2\). Call such \(A\) accepted in \(C\).
The rule in \(C\) says that the accepted events are closed under binary
intersection, since \((A\cap B)^c=A^c\cup B^c\). By induction, they are then
closed under every finite nonempty intersection. Let \(C^+\) be the set of
positive-mass states in \(C\). Only \(C^+\) contributes to \(P(A\cap C)\).

Sufficiency. Assume the displayed condition. If \(C^+=\{s\}\), the accepted
events are those containing \(s\). If \(C^+=\{s,t\}\) with \(p_s>p_t\), then
\(p_s>T/2\), and an event missing \(s\) has \(P(A\cap C)\leq p_t<T/2\), so the
accepted events are again those containing \(s\). If \(p_s=p_t\), the accepted
events are those containing both \(s\) and \(t\). If \(C^+\) has at least three
states, let \(s_i\) be its member of least index. The other members have
indices above \(i\), and at least two of them exist, so \(i\leq k-2\). Then
\[
 p_i>\sum_{j=i+1}^kp_j\geq T-p_i ,
\]
so \(p_i>T/2\). An event missing \(s_i\) has \(P(A\cap C)\leq T-p_i<T/2\).
Hence the accepted events are exactly those containing \(s_i\). In every
case the accepted family is the set of events containing a fixed nonempty
set of states, and it is closed under intersection.

Necessity. Suppose the condition fails at some \(i\leq k-2\), so
\(p_i\leq\sum_{j>i}p_j\). Let \(C=\{s_i,\ldots,s_k\}\), with \(T=P(C)\). Then
\(p_i\leq T/2\). For each \(j>i\), \(p_j\leq p_i\leq T/2\). Equality
\(p_j=T/2\) would force \(p_i=p_j=T/2\), which leaves no mass for the third
positive state of \(C\); so \(p_j<T/2\). Hence \(C\setminus\{s_j\}\) is
accepted in \(C\), since its mass is \(T-p_j>T/2\). There are at least two
such \(j\). The intersection of all these events is \(\{s_i\}\), and its mass in \(C\) is \(p_i\leq T/2\), so it is not
accepted. This contradicts closure under finite intersection.

For the two examples, the uniform law on two states has \(k=2\), so the
condition is vacuous and the rule holds. On three states, \(i=1\) requires
\(1/3>2/3\), which is false, so the rule fails. Concretely, in \(C=S\) each
two-element event has mass \(2/3>1/2\), while the intersection of two of them
has mass \(1/3\).
\end{proof}\end{annexitem}
